\section{Related Work}
\label{sec:related}

\noindent\textbf{Reinforcement Learning for LLM Reasoning.}
RL now complements supervised fine-tuning (SFT) to enhance LLM reasoning.
Instability and cost in PPO have motivated alternatives such as Group Relative Policy Optimization (GRPO)~\cite{shao2024grpo,li2025mogrp} and Direct Preference Optimization (DPO)~\cite{goffer2023dpo}.
Typical pipelines start from an SFT checkpoint with Chain-of-Thought traces and apply GRPO (or derivatives) to improve factual consistency and deductive depth~\cite{yue2025cot,zhang2025sgrpo}, often incorporating RL from Human Feedback (RLHF)~\cite{outrun2023rlhf} or Reinforcement Learning from AI Feedback (RLAIF)~\cite{rajesh2024rlaif}.
Multi-objective settings introduce label-specific rewards or curricula that adapt weights during training~\cite{chen2025emorl,shibata2025dynamicrl}, and recent work has explored hierarchical reward decomposition to scale reasoning depth~\cite{nakano2023hapt}.
PhyPrompt follows this two-stage recipe: SFT on a physics-focused CoT corpus, then GRPO with a curriculum reward that gradually shifts emphasis from semantic alignment to physical-commonsense cues vital for video generation.

\noindent\textbf{Benchmarks for Physical Commonsense in Video Generation.}
To judge physics realism, purpose-built datasets replace pixel-level metrics.
VideoPhy2 provides human-plus-automatic scores for Semantic Adherence (SA) and Physical Commonsense (PC)~\cite{bansal2025videophy2}, while PhyGenBench widens coverage to 160 prompts across 27 physical laws~\cite{phygenbench2024}.
PhyCoBench tracks coherence in seven principle categories via optical-flow signals~\cite{chen2025phycobench}, and IPV-Bench stresses models with intentionally impossible scenes~\cite{wang2025ipvbench}.
Other recent benchmarks include VisualPhysicist for dynamic object interactions~\cite{singh2024visualphysicist} and SpeedNet for temporal consistency in motion~\cite{kim2023speednet}.
These resources form the first quantitative yardsticks for physics-aware video synthesis.

\noindent\textbf{Prompt Rewriting for Image and Video Generation.}
Promptist uses RL-based adaptation to refine text-to-image prompts~\cite{hao2023promptist}, and DiffusionPrompt employs gradient-based style rewrites~\cite{lee2023diffusionprompt}. In text-to-video, VPO combines SFT and DPO with multi-level feedback~\cite{cheng2025vpo}, Prompt-A-Video applies reward-driven caption evolution~\cite{ji2024prompt}, MotionPrompt adds optical-flow temporal guidance~\cite{nam2024motionprompt}, MotionCraft and ReVision inject physics cues at generation~\cite{agnolucci2024motioncraft,wang2025revision}, and newer approaches, VideoPromptTuning~\cite{xu2024vpt}, Prompt2Video~\cite{fernandez2024prompt2video}, and PhyT2V~\cite{xue2024phyt2v}, use lightweight adapters or iterative CoT reasoning for enhanced physical adherence.

PhyPrompt differs by pre-training on a physics-focused CoT corpus and maximising an automatic PC reward computed on the generated video, yielding a lightweight, generator-agnostic front end for physics-aware T2V models.
Besides, PhyPrompt is the first prompt rewriter that trains end-to-end with RL on video-based physical-commonsense rewards and transfers zero-shot across multiple heterogeneous T2V generators without re-tuning.
These qualities position PhyPrompt as a practical bridge between user intent and physics-consistent video generation.
